% sec3.tex — c03｜原文 p13（§3 标题下的正文起）–p19，并含 p20 上半至 §4 标题前
% §3 超对称量子力学与 K 理论（§3.1–§3.6）；本块公式号 (21)–(38)，共 18 条；编号环境 3.1–3.26
% 块边界：式 (20) 在 p13 顶部、§3 标题之上的 §2 例 2.22 内，不属本块。
% 图 3（p15 顶部）、图 4（p18 顶部）一律从略，用 \figomit。
% [Tutor Note: 底本 p13 页首的式 (20) 排在 §3 标题之前，属 §2 例 2.22 中 A-扭场论的数据；故本节从式 (21) 起，式 (20) 见第 2 节末。]
% [Tutor Note: 注 3.25 末句与例 3.26 排在 p20 页首、§4 标题之前，编号仍属 §3，故本块译至 §4 标题前止；这是底本的分页所致，非编号错乱。]

% [原文 p.13]
本节的目标是解释猜想 1.1 的下述类比。

\begin{conjecture}[3.1（Stolz–Teichner）]
对每个 $n\in\Z$，存在虚线箭头
\begin{equation}
\begin{tikzpicture}[baseline=(current bounding box.center), >=stealth]
  \node (A) at (0,3.4) {$\left\{\begin{array}{c}\text{一族 $n$ 维 spin$^c$}\\[2pt] \text{流形 $M\to X$}\end{array}\right\}$};
  \node (C) at (9.4,1.7) {$\mathrm{K}^{-n}(X)$};
  \node (B) at (0,0) {$\left\{\begin{array}{c}\text{$X$ 上次数为 $-n$ 的}\\[2pt] \text{1 维超对称}\\[2pt] \text{场论}\end{array}\right\}$};
  \draw[->] (A.south east) -- node[above, sloped] {拓扑指标} (C.north west);
  \draw[->, dashed] (A.south) -- node[left] {量子化} (B.north);
  \draw[->, dashed] (B.east) -- node[below right=1pt, pos=0.42] {上闭链${}_{-n}$} (C.south west);
\end{tikzpicture}
\tag{21}
\end{equation}
使三角形交换；其中“拓扑指标”是 Atiyah–Bott–Shapiro 的 spin$^c$ 定向 [ABS64]，“量子化”是对 $M$ 上超对称力学的量子化，“上闭链$_n$”把 $X$ 上的一个量子理论送到 $\mathrm{K}^n(X)$ 中的一个类。当 $n$ 为负时，空流形是唯一的 spin 流形，此时 (21) 的内容是存在一个上闭链映射。
\end{conjecture}
% [Tutor Note: 猜想 3.1 的流形是 spin^c 流形（式 (21) 图中为 spin^c），而正文此处说 "the only spin manifold"；照译，应读作 spin^c 流形。]
% [Tutor Note: 式 (21) 的箭头标签是 cocycle_{−n}，紧邻正文却写 cocycle_n；与式 (1) 同一情形，下标随度数变号，二者不矛盾。]

除某些技术性告诫外，猜想 3.1 是一个定理。事实上，它已对 [ST04, §3.2] 与 [HST10] 中 1 维超对称场论的定义得到验证，其综述见 §3.3 与下文假设 3.20。然而，人们相信这些既有论证难以很好地推广到 $2|1$ 维情形，因此猜想 3.1 的各种变体仍然悬而未决；其中一个变体见注 3.24。

% [原文 p.14]
\subsection{超对称量子力学的动机}
给定定义 2.1 意义下的一个量子力学系统 $(\cH,H)$，设 $\cH=\cH^+\oplus\cH^-$ 带有一个 $\Z/2$-分次，且 $H$ 尊重该分次（即 $H$ 是偶映射）。等价地，$\cH$ 上带有一个对合 $(-1)^F\colon\cH\to\cH$，其特征子空间为
\begin{equation}
(-1)^F=\begin{cases}+1 & \text{在 } \cH^+ \text{ 上}\\ -1 & \text{在 } \cH^- \text{ 上}\end{cases}
\qquad\text{且}\qquad H(-1)^F=(-1)^F H.
\tag{22}
\end{equation}
超对称量子力学在这些数据之上增添 $\cN$ 个奇算子 $Q_i$，满足
\begin{equation}
Q_i^2=H,\quad [Q_i,Q_j]=Q_iQ_j+Q_jQ_i=0\ (i\ne j),\quad [(-1)^F,Q_i]=0,\quad i\in\{1,2,\dots,\cN\}
\tag{23}
\end{equation}
其中 $[-,-]$ 是分次交换子 [Wit82a, §2]。$Q_i$ 是超对称算子。

\begin{example}[3.2（定向流形上的 $\cN=2$ 超对称量子力学）]
超对称量子力学的经典例子取状态空间为 $\cH=\Omega^\bullet(M)$，即定向 Riemann 流形 $M$ 上的微分形式 [Wit82b, §2]。此时哈密顿量定义为 Hodge 拉普拉斯算子 $H=\dd\dd^*+\dd^*\dd$，它有一对交换的奇平方根
\begin{equation}
Q_1=\dd+\dd^*,\qquad Q_2=\sqrt{-1}(\dd-\dd^*),\qquad [Q_1,Q_2]=0,\qquad Q_1^2=Q_2^2=H.
\tag{24}
\end{equation}
这是 $\cN=2$ 超对称的数据。忘掉 (24) 中一个 $Q_i$ 便抽取出一个 $\cN=1$ 理论。其他 $\cN=1$ 理论来自（扭的）Dirac 算子，见例 3.13。
\end{example}

猜想 3.1（以及下文）的焦点是 $\cN=1$ 超对称。此时，(23) 中的算子确定超 Lie 群 $\R^{1|1}$ 的超 Lie 代数的一个表示，$\R^{1|1}$ 的乘法为\footnote{这个公式应在 $\R^{1|1}$ 的函子点（functor of points）意义下理解，例如参见 [DM99]。}
\begin{equation}
(t,\theta)\cdot(s,\eta)=(t+s+\theta\eta,\theta+\eta),\qquad (t,\theta),(s,\eta)\in\R^{1|1}.
\tag{25}
\end{equation}
沿 $\theta=0$ 包含，得到通常加法下的 Abel 子群 $\R<\R^{1|1}$。沿此包含作延拓，就把通常的时间演化 (6) 推广为半超群（supersemigroup）的一个表示 [HST10, §5] [ST04, §3.2]
\begin{equation}
\R^{1|1}_{\ge0}\to\mathrm{End}(\cH),\qquad (t,\theta)\mapsto\ee^{-tH+\theta Q},
\tag{26}
\end{equation}
其中 $(t,\theta)$ 是子超流形 $\R^{1|1}_{\ge0}\subset\R^{1|1}$（$t\ge0$）上的坐标。
% [Tutor Note: p14 的两条脚注在原文中为第 5 号（$\R^{1|1}$ 的函子点，见上）与第 6 号（奇无界算子的定义域，见 §3.2）；译本由 LaTeX 自动编号，号码不必相同。]

\begin{definition}[3.3]
具有（$\cN=1$）超对称的一个量子力学系统是一对 $(\cH,Q)$，其中 $Q$ 是 Hilbert 空间 $\cH$ 上的奇无界自伴算子，它生成一个迹类（trace class）的超半群表示 (26)。若存在酉映射 $u\colon\cH\to\cH'$ 使 $Q'=uQu^{-1}$，则称系统 $(\cH,Q)$ 与 $(\cH',Q')$ 等价。
\end{definition}

我们在下文展开这一点。

\subsection{具有 $\cN=1$ 超对称的次数为 $n$ 的量子系统的空间}
与 §2.1 中描述的哈密顿量空间类似，我们把定义 3.3 推广为带超对称的量子系统的空间——并最终推广为一个栈。着眼 K 理论上的 $\Z$-分次，我们还纳入 Clifford 模结构；在物理上，这些结构来自自由费米子的（射影）对称性代数，见例 2.11 与注 2.23。

固定一个无穷维可分 $\Z/2$-分次 Hilbert 空间 $\cH$，其偶子空间与奇子空间都是无穷维的。置 $\cH_n:=\mathrm{Cl}_n\otimes\cH$，其中 $\mathrm{Cl}_n$ 是第 $n$ 个 Clifford 代数。设 $Q$ 是 $\cH_n$ 上的 Clifford 线性奇自伴无界算子。\footnote{$\Z/2$-分次 Hilbert 空间上的奇无界算子，其定义域 $\mathrm{dom}(Q)\subset\cH$ 是分次子空间，且 $Q\colon\mathrm{dom}(Q)\to\cH$ 是奇线性映射。} 注意 $H=Q^2$ 是 Clifford 线性的、非负的偶算子。对 $H=Q^2$ 应用谱定理，得到一族互相正交的特征子空间 $\{\cH_\lambda\subset\cH_n\}_{\lambda\ge0\in\R}$，如
% [原文 p.15]
图 1 所示。\figomit{3}{推广图 1：超对称把标号 $\cH_\lambda\subset\cH$ 增强为一个 $\Z/2$-分次以及一个满足 $Q_\lambda^2=\lambda\cdot1$（$\lambda\ge0$）的奇算子 $Q_\lambda$。沿用 [Wit82a]，我们用一个 $\circ$ 记一维偶空间、用一个 $\times$ 记一维奇空间（分别代表玻色态与费米态）。由于 $\lambda\ne0$ 时 $Q_\lambda$ 可逆，$\cH_\lambda$（$\lambda\ne0$）上的偶维数必须等于奇维数。因此，零特征子空间中偶维数与奇维数之差是形变不变量 $[\cH_0]\in\mathrm{K}(\mathrm{pt})$。}
这些子空间是有限秩 $\Z/2$-分次 $\mathrm{Cl}_n$-模，带有奇算子 $Q_\lambda$ 的作用，其特征值为 $\pm\sqrt{\lambda}$。当 $\lambda\ne0$ 时，算子 $\frac{1}{\sqrt{\lambda}}Q$ 在 $\cH_\lambda$ 上生成一个与 $\mathrm{Cl}_n$-作用交换的 $\mathrm{Cl}_{-1}$-作用。因此，$\lambda\ne0$ 时的 $\cH_\lambda$ 获得 $\mathrm{Cl}_{-1}\otimes\mathrm{Cl}_n$-模的结构。于是这些构形的形变（例如见图 3）重新得到 Atiyah–Bott–Shapiro 对 $\pi_*\mathrm{K}^n$ 的系数的描述：$\mathrm{Cl}_n$-模的同构类模掉那些可延拓为 $\mathrm{Cl}_{-1}\otimes\mathrm{Cl}_n$-模者 [ABS64]。
% [Tutor Note: 原文把指标集写作 $\{\cH_\lambda\subset\cH_n\}_{\lambda\ge0\in\R}$；照排，应读作 $\lambda\in\R_{\ge0}$。]
% [Tutor Note: 图 3 图注中的 "o's" 指图上小圆圈 $\circ$；本译作「一个 $\circ$ 记一维偶空间」，与图中记号一致。]

\begin{definition}[3.4]
记 $\sHam^n$ 为生成迹类超半群表示 (26) 的奇自伴 $\mathrm{Cl}_n$-线性无界算子 $Q$ 构成的空间。
\end{definition}

\begin{remark}[3.5]
无界算子 $Q$ 何时确定一个迹类超半群表示，其分析条件见 [HST10, §5]。
\end{remark}
% [Tutor Note: 定义 3.4 与定义 3.7 在原文中用上标 $\sHam^n$、$\SQM^n$，而命题 3.12 的证明与式 (31) 等处用下标；照原文混用，不改。]

\begin{definition}[3.6]
对 $\Z/2$-分次 Hilbert 空间，定义（超）酉群
\[
\mathrm{U}^\pm(\cH)\simeq\left\{u\in\mathrm{U}(\cH)\ \middle|\ u=\begin{bmatrix}*&0\\0&*\end{bmatrix}\ \text{或}\ \begin{bmatrix}0&*\\*&0\end{bmatrix}\right\}
\]
其中 $u$ 相对分解 $\cH=\cH^+\oplus\cH^-$（偶子空间与奇子空间）取分块形式。群 $\mathrm{U}^\pm(\cH)$ 是分次的（例如在 [Fre12] 的意义下），分次由同态
\[
\mathrm{U}^\pm(\cH)\to\pi_0(\mathrm{U}^\pm(\cH))\simeq\Z/2
\]
给出，后者读出酉映射保持还是反转 $\cH$ 上的分次。记 $\mathrm{PU}^\pm(\cH):=\mathrm{U}^\pm(\cH)/\mathrm{U}(1)$ 为射影化。类似地，对分次酉 $\mathrm{Cl}_n$-模 $\cH_n$，定义 $\mathrm{U}^\pm(\cH_n)$ 为与 Clifford 作用交换的酉变换，并记 $\mathrm{PU}^\pm(\cH_n)=\mathrm{U}^\pm(\cH_n)/\mathrm{U}(1)$ 为射影化。
\end{definition}

\begin{definition}[3.7]
定义拓扑群胚 $\SQM^n:=\sHam^n/\!/\mathrm{PU}^\pm(\cH_n)$ 为酉群的共轭作用所诱导出的作用的商。
\end{definition}

我们有定义 2.8 的下述超对称推广。

\begin{definition}[3.8]
设 $G$ 是紧 Lie 群。带 $G$-对称性的次数 $n$ 的超对称量子力学系统构成的栈，是映射栈 $\mathrm{Map}(\mathrm{pt}/\!/G,\SQM^n)$。
\end{definition}
展开之，$\mathrm{Map}(\mathrm{pt}/\!/G,\SQM^n)$ 中的一个点是同态 $\alpha\colon G\to\mathrm{PU}^\pm(\cH)$，它定义 $G$ 的一个中心扩张 $G^\alpha$ 以及一个酉表示 $G^\alpha\to\mathrm{U}^\pm(\cH)$（类似于 (8)），连同 $\cH$ 上与 $G^\alpha$-作用交换的一个超对称生成元 $Q$。复合 $G^\alpha\to\mathrm{U}^\pm(\cH)\to\Z/2$ 赋予群 $G^\alpha$ 一个分次，它读出哪些元素以保分次或反分次的自同构作用在 $\cH$ 上。

% [原文 p.16]
\subsection{超对称量子力学与 K 理论}
Atiyah–Bott–Shapiro 对分次环 $\pi_{-*}\mathrm{K}=\mathrm{K}^*(\mathrm{pt})$ 的描述 [ABS64] 给出一个映射
\begin{equation}
\sHam^n\to\mathrm{K}^n(\mathrm{pt}),\qquad Q\mapsto[\mathrm{Index}(Q)]:=[\mathrm{Ker}(Q)]=[\cH_0]
\tag{27}
\end{equation}
它把超对称生成元 $Q$ 送到其核，视作一个有限秩 $\mathrm{Cl}_n$-模。

\begin{definition}[3.9]
$Q\in\sHam^n$ 的 Witten 指标是它在 (27) 下的像。
\end{definition}

事实上，我们得到下述结论。

\begin{theorem}[{3.10（[HST10]，[Che08]）}]
空间 $\sHam^n$ 表示函子 $\mathrm{K}^n(-)$。
\end{theorem}
\begin{proof}[证明梗概]
空间 $\sHam^n$ 同伦等价于一个 Clifford 线性 Fredholm 算子空间，后者已知表示 K 理论谱中的第 $n$ 个空间 [AS69]。事实上，$Q\in\sHam^n$ 中算子特征值的增长条件保证：由泛函演算定义的算子 $Q/(1+Q^2)$ 是 Fredholm 的。
\end{proof}

\begin{example}[3.11（拓扑理论）]
推广例 2.5：任何有限维 $\Z/2$-分次子模 $V\subset\cH_n$ 都通过取 $Q=0$、定义域为 $V$ 而确定 $\sHam^n$ 中的一个点。这对应于支撑在 $0\in\R_{\ge0}$、标号为 $\cH_0=V$ 的构形。我们注意到，自然包含
\begin{equation}
\left\{\begin{array}{c}\text{次数 $n$}\\[2pt]\text{拓扑理论}\end{array}\right\}\xrightarrow{\ \text{群完备化}\ }\sHam^n\to\SQM^n
\tag{28}
\end{equation}
可以与 $\mathrm{Cl}_n$-模的群完备化等同起来。
\end{example}

\begin{proposition}[3.12]
$\{\SQM^n\}_{n\in\Z}$ 表示拓扑栈的扭 K 理论，特别地表示紧 Lie 群 $G$ 的 $G$-空间上的扭等变 K 理论 [FHT11a]。
\end{proposition}
\begin{proof}[证明略述]
给定拓扑栈的一个映射 $\alpha$，
\begin{equation}
\mathfrak X\xrightarrow{\ \alpha\ }\SQM^n=\sHam^n/\!/\mathrm{PU}^\pm(\cH_n)
\tag{29}
\end{equation}
与投影作后复合，便确定一个映射 $\tau\colon\mathfrak X\to\mathrm{pt}/\!/\mathrm{PU}^\pm(\cH_n)$，它分类一个带 Clifford 作用的射影 Hilbert 丛 $\PP\cH^\tau\to\mathfrak X$（也称为扭 Hilbert 丛 [FHT11a, §3.2]）。此外，该丛由分次 gerbe $\mathrm{pt}/\!/\mathrm{U}^\pm(\cH_n)\to\mathrm{pt}/\!/\mathrm{PU}^\pm(\cH_n)$ 沿 $\tau$ 的拉回来分类。映射 $\alpha$ 的其余数据，是 $\mathfrak X$ 上以 $\sHam_n$ 为纤维的配丛的一个截面。由于 $\sHam_n$ 表示 $\mathrm{K}^n(-)$，该截面的同伦类是一个扭 K 理论类。特殊化到（好的）空间 $X$ 在紧 Lie 群 $G$ 作用下的商栈 $\mathfrak X=X/\!/G$，便重新得到 $X$ 的 $G$-等变扭 K 理论 [FHT11a, §3.5.4]。
\end{proof}

\subsection{超对称量子力学与指标定理}

\begin{example}[3.13（spin$^c$ 流形上的超对称量子力学）]
设 $M$ 是 Riemann 流形，考虑路径 $x\colon\R\to M$ 连同（费米子）$\psi\in\Gamma(\R;x^*TM)$ 的经典理论，其作用量为
\begin{equation}
S(x,\psi)=\frac12\int_\R(\|\dot x\|+\langle\psi,\nabla_{\dot x}\psi\rangle)\dd t.
\tag{30}
\end{equation}
设 $M$ 带有一个选定的 spin$^c$ 结构，其旋量丛 $\mathrm{S}:=P\times_{\mathrm{Spin}(n)}\mathrm{Cl}_n$ 支撑 $\mathrm{Cl}_{-n}$-线性的 Dirac 算子 $\mathord{\not{D}}$。\footnote{右 $\mathrm{Cl}_n$-模之丛 $\mathrm{S}=P\times_{\mathrm{Spin}(n)}\mathrm{Cl}_n$ 确定一个左 $\mathrm{Cl}_{-n}=\mathrm{Cl}^{op}_n$-模之丛。} 那么经典理论 (30) 量子化后状态空间为 $\cH=\Gamma(\mathrm{S})$、哈密顿量为 $H=\mathord{\not{D}}^2$ [Fre99, 第 43–44 页]。取 $Q=\mathord{\not{D}}$ 便按定义 3.7 的意义定义一个次数为 $-n$ 的超对称量子系统，从而给出
\begin{equation}
\left\{\begin{array}{c}\text{$n$ 维 spin$^c$}\\[2pt]\text{流形 $M$}\end{array}\right\}\to\SQM_{-n},\qquad M\mapsto\mathord{\not{D}}.
\tag{31}
\end{equation}
若 $M$ 带保持 spin$^c$ 结构的 $G$-作用，则 $\Gamma(\mathrm{S})$ 自然地是 $G$-表示。若改为 $W_3(TM)\in\mathrm{H}^3_G(M;\Z)$ 从 $BG$ 的上同调（即点的等变上同调）拉回，则 $\Gamma(\mathrm{S})$ 只有射影的 $G$-作用。两种情形都给出定义 3.8 意义下带 $G$-对称性的超对称量子力学系统。
\end{example}
% [Tutor Note: 式 (30) 的动能项原文作 $\|\dot x\|$（未平方）；通常写作 $\|\dot x\|^2$。照原文排，疑为笔误，已登记 ERRATA。]
% [Tutor Note: p16 的脚注为原文第 7 号（右 $\mathrm{Cl}_n$-模丛与左 $\mathrm{Cl}_{-n}$-模丛的对应），译本自动编号。]

% [原文 p.17]
\begin{proof}[猜想 3.1 的证明梗概]
若把 $X$ 上次数为 $-n$ 的 1 维超对称欧几里得场论定义为映射 $X\to\Ham(\cH_{-n})$（比较 [ST04, Definition 3.2.10]），则 (21) 中的上闭链映射由命题 3.12 得出。给定 spin$^c$ 流形的一族 $M\to X$，逐纤维的 $\mathrm{Cl}_{-n}$-线性 Dirac 算子给出 (31) 的族版本，从而产生一个映射 $X\to\Ham(\cH_{-n})$，它对应于族解析指标 [ST04, Remark 3.2.21]。三角形 (21) 的交换性就是族指标定理 [AS71]。
\end{proof}

\begin{remark}[3.14]
指标定理的物理论证考察 $M$ 上超对称量子力学的时间演化半群的极限
\begin{equation}
\mathrm{sTr}_{\mathrm{Cl}_{-n}}\bigl(\ee^{-t\mathord{\not{D}}^2}\bigr)=\begin{cases}\mathrm{Index}(\mathord{\not{D}}) & t\to\infty\\[6pt] \displaystyle\int_X\hat{A}(X) & t\to0.\end{cases}
\tag{32}
\end{equation}
McKean–Singer 定理 [MS67] 表明该超迹不依赖于 $t$，并把 $t\to\infty$ 极限与 $Q=\mathord{\not{D}}$ 的（Witten）指标等同起来。另一方面，$t\to0$ 极限可用路径积分技巧计算。具体地，超对称局部化把该路径积分等同于 $\hat{A}$-形式在 $M$ 上的积分 [Wit82b, AG83]。数学处理参见 [Ati85, JP90, BGV92]。
\end{remark}

K 理论中的拓扑前推来自 Thom 同构，后者也有一个用超对称量子力学的描述。

\begin{example}[3.15（自由费米子的超对称形变）]
由例 3.11，例 2.11 与 2.24 中的自由费米子有到次数 $n$ 超对称理论的一个典范延拓：状态空间是分次 $\mathrm{Cl}_n$-模 $V$，超对称生成元为 $Q=0$。考虑修改后的自由费米子经典场论 (9)
\begin{equation}
\psi\colon\R\to\R^n,\qquad S(\psi)=\frac12\int_\R(\langle\psi,\partial_t\psi\rangle+\langle x,\psi\rangle)\dd t
\tag{33}
\end{equation}
它把源项 $\langle x,\psi\rangle$ 加到作用量上，其中 $x\in\R^n$ 固定。该系统的量子化给出状态空间 $V=\mathrm{Cl}_n$（作为左 $\mathrm{Cl}_n$-模），连同由超对称生成元 $Q=\mathrm{cl}_x\colon V\to V$ 决定的时间演化算子，其中 $\mathrm{cl}_x$ 是沿 $x\in\Pi\R^n\subset\mathrm{Cl}_n$ 的右 Clifford 乘法。让 $x\in\R^n$ 变动，得到映射
\begin{equation}
\R^n\to\SQM^n,\qquad x\mapsto Q=\mathrm{cl}_x
\tag{34}
\end{equation}
它在 $x=0$ 处是拓扑理论。可以把 (34) 与 K 理论中第 $n$ 个悬置类等同起来 [ABS64]。存在 (33) 上的 $\mathrm{O}_n$-作用，它旋转 $\psi$ 与源项 $x$，而映射 (34) 是（射影地）$\mathrm{O}_n$-等变的。这给出 (10) 的推广
\begin{equation}
\Th_n\colon\R^n/\!/\mathrm{O}_n\to\SQM^n=\sHam^n/\!/\mathrm{PU}^\pm(\cH),\qquad
\begin{gathered}\mathrm{O}_n\to\mathrm{PU}^\pm(\cH),\\ x\mapsto\mathrm{cl}_x=Q(x),\quad x\in\R^n\end{gathered}
\tag{35}
\end{equation}
其中映射 $\mathrm{O}_n\to\mathrm{PU}^\pm(\cH_n)$ 分类 $\mathrm{Pin}^c$-扩张。在命题 3.12 之下，(35) 底层的扭 K 理论类是第 $n$ 个泛（扭等变）Thom 类 $[\Th_n]\in\mathrm{K}^{w_1+W_3}(\R^n/\!/\mathrm{O}_n)_c\simeq\mathrm{K}^{w_1+W_3}_{\mathrm{O}_n}(\R^n)_c$ [FHT11a, §3.6] [ST04, Remark 3.2.22]。
\end{example}
% [Tutor Note: 式 (35) 中写 $\mathrm{PU}^\pm(\cH)$，紧接正文写 $\mathrm{PU}^\pm(\cH_n)$；原文如此，照排。]
% [Tutor Note: Thom 类的次数写作 $w_1+W_3$（小写 $w_1$、大写 $W_3$）；通常记作 $w_1+w_3$，$W_3$ 是 spin^c 结构的整提升。照原文排。]

\begin{remark}[3.16]
Mathai–Quillen 形式化 [MQ86] 计算 (33) 的配分函数，重新得到由 $\mathfrak{o}_n$ 上（逆）$\hat{A}$-形式所决定的一个泛 Thom 类。
\end{remark}
% [Tutor Note: 上句原文末为 "on o_n"，其中 o_n 是 Fraktur 体 $\mathfrak{o}_n$（$\mathrm{O}_n$ 的 Lie 代数），文字层易误读作 "on on"。]

\subsection{配边框架中的超对称量子力学}
可以把超对称量子力学重新打包成 $1|1$-欧几里得配边范畴 $\EBord_{1|1}$ 的一个表示，它推广 $\EBord_1$。$\EBord_{1|1}$ 的对象是 $\bullet=\mathrm{spt}^+$ 与 $\circ=\mathrm{spt}^-$ 的有限不交并，即带两种定向选取的超点。态射见图 4，其中 $(t,\theta)\in\R^{1|1}_{\ge0}$ 编码一个“超长度”。与之前类似，我们得到与配边复合相关联的关系，例如
\begin{equation}
\mathrm{I}^\pm_{s,\eta}\circ\mathrm{I}^\pm_{t,\theta}\simeq\mathrm{I}^\pm_{s+t\pm\eta\theta,\eta+\theta},
\tag{36}
\end{equation}
它们给出 (26) 中超半群 $\R^{1|1}_{\ge0}$ 中乘法的几何对应物。

\figomit{4}{$\EBord_{1|1}$ 中的态射是上述生成元的有限不交并。它们把图 2 中的生成元加以增强：加入 spin 结构，并把长度参数 $t\in\R_{\ge0}$ 换成取值于 $\R^{1|1}_{\ge0}$ 的“超长度”$(t,\theta)$。}

% [原文 p.18]
\begin{remark}[3.17]
维数为 $1|1$ 的超欧几里得几何由超 Lie 群 (25) 按如下意义决定：人们规定群 $\R^{1|1}$ 以欧几里得等距的方式作用在超流形 $\R^{1|1}$ 上。这类似于维数为 $1$ 的欧几里得几何由下述规定决定：群 $\R$ 是 $\R$ 的保向对称群，通过平移作用。此外还有一个由 $\theta\mapsto-\theta$ 确定的 $1|1$-欧几里得等距，起 spin 翻转的作用。因此，$1|1$-欧几里得场论推广了例 2.24 中的 spin 欧几里得场论。下文我们隐式地施加相应的 $B\Z/2$-等变性质，细节参见 [ST11, §4.3]。
\end{remark}

沿用 (19)，我们对每个 Clifford 代数 $\mathrm{Cl}_n$ 构造一个扭。

\begin{definition}[3.18（次数扭）]
次数扭由下述指派决定：
\begin{align}
&\alpha_n\colon\EBord_{1|1}\to\Alg,\notag\\
&\alpha_n(\mathrm{spt}^+)=\mathrm{Cl}_n,\quad \alpha_n(\mathrm{spt}^-)=\mathrm{Cl}^{op}_n=\mathrm{Cl}_{-n},\notag\\
&\alpha_n(\mathrm{I}^\pm_{t,\theta})={}_{\mathrm{Cl}_n}(\mathrm{Cl}_n)_{\mathrm{Cl}_n},\quad
 \alpha_n(\cap^+_{t,\theta})=(\mathrm{Cl}_n)_{\mathrm{Cl}_n\otimes\mathrm{Cl}^{op}_n},\notag\\
&\alpha_n(S^\pm_{t,\theta})=\mathrm{Cl}_n\otimes_{\mathrm{Cl}_n\otimes\mathrm{Cl}^{op}_n}\mathrm{Cl}_n,\notag\\
&\alpha_n(\cup^+_{t,\theta})={}_{\mathrm{Cl}_n\otimes\mathrm{Cl}^{op}_n}\mathrm{Cl}_n,\tag{37}
\end{align}
其反射结构来自 Clifford 代数上的 $\ast$-结构 $\ast\colon\mathrm{Cl}_n\xrightarrow{\sim}\ol{\mathrm{Cl}^{op}_n}=\ol{\mathrm{Cl}}_{-n}$。这些值与 $(t,\theta)$ 无关，因此次数扭是拓扑的。次数 $n$ 的扭场论的取值可以类比 (20) 读出。Morita 等价 $\mathrm{Cl}_n\otimes\mathrm{Cl}_m\simeq\mathrm{Cl}_{n+m}$ 给出扭的同构 $\alpha_n\otimes\alpha_m\simeq\alpha_{n+m}$（$n,m\in\Z$）。
\end{definition}
% [Tutor Note: 式 (37) 的 ${}_{\mathrm{Cl}_n}(\mathrm{Cl}_n)_{\mathrm{Cl}_n}$ 是双模记号：左作用写在左、右作用写在右；$\mathrm{Cl}^{op}_n=\mathrm{Cl}_{-n}$ 由 Clifford 代数的 $8$ 周期性与取反代数给出。]
% [Tutor Note: 式 (37) 中 $\otimes_{\mathrm{Cl}_n\otimes\mathrm{Cl}^{op}_n}$ 是平衡张量积（coequalizer），不是分次张量积。]

\begin{remark}[3.19]
Clifford 代数都是 2-范畴 $\Alg$ 中的可逆对象，因此是全对偶的（fully dualizable）。由配边假设 [Lur09a]，$\mathrm{Cl}_n$ 确定一个全扩张的 2 维拓扑场论。该理论在 $0$ 维与 $1$ 维流形上的取值由 (37) 给出 [Gun16]。因此，次数扭是这个全扩张拓扑场论的一个截断，其反射结构继承自 $\mathrm{Cl}_n$ 上的 $\ast$-结构。
\end{remark}
% [Tutor Note: 原文 "with reflection structure inherited the ∗-structure" 缺介词 from，应为 inherited from；中文译文已按正确意思作「继承自」，已登记 ERRATA。]

令 $\QFT^\alpha_{1|1}(\mathfrak X)$ 记 $\mathfrak X$ 上 $\alpha$-扭的反射正 $1|1$ 维欧几里得场论的群胚（见定义 1.9）。反射正性与定义 2.13 中的基本相同。我们先分析次数扭 (37) 对应的群胚 $\QFT^n_{1|1}:=\QFT^{\alpha_n}_{1|1}(\mathrm{pt})$。

\begin{hypothesis}[{3.20（比较 [ST04, Proposition 3.2.6] 与 [HST10, Corollary 6.38]）}]
存在栈的满射（epimorphism）$\QFT^n_{1|1}\to\SQM^n$。
\end{hypothesis}
\begin{proof}[略证]
论证与假设 2.14 的论证类似，因此我们将（更加）简略。函子 $\QFT^n_{1|1}\to\SQM^n$ 把 $1|1$-欧几里得场论 $E$ 送到其半超群表示 $E(\mathrm{I}^+_{t,\theta})$ 的生成元 $Q$。确实，由 [HST10, Lemma 5.10]，有 $E(\mathrm{I}^+_{t,\theta})=\ee^{-tH+\theta Q}$，其中 $Q$ 是由此唯一确定的无界奇 Clifford 线性算子，定义域为 $E(\mathrm{spt}^+)\subset\cH$。

为说明该指派确定一个满射，$\EBord_{1|1}$ 的生成元与关系表示（比较 [HST10, Theorem 6.42]）表明：次数 $n$ 的 $1|1$-欧几里得场论由自伴 Clifford 模 $\bigl(E(\mathrm{spt}^+),\langle-,-\rangle\bigr)$ 以及来自 $E(\mathrm{I}^+_{t,\theta})$ 的自伴 $\mathrm{Cl}_n$-线性半超群表示决定。类比于 (11) 的构造随后确定所要的单侧逆函子。
\end{proof}

% [原文 p.19]
\begin{remark}[3.21]
对栈的等价 $\QFT^n_{1|1}\to\SQM^n$ 的期望会遇到一个障碍：存在许多互不同构的拓扑向量空间 $E(\mathrm{pt}^+)$，它们导致互不同构的场论，但这些场论在 $\SQM^n$ 中作为对象却是同构的（在把一个拓扑内积空间完备化为 Hilbert 空间之后）。一种可能的修补是施加一个更弱的等价关系。场论的 concordance 类 [ST11, Definition 1.12] 是最有希望的选择，并且人们期望：$X$ 上场论的 concordance 类与从 $X$ 到栈 $\QFT^n_{1|1}$ 的映射的 concordance 类一一对应。然而，concordance 等价 $\QFT^n_{1|1}\to\SQM^n$ 仍要求对场论的定义作进一步的技术性修改，参见 [ST11, Remark 3.16] 与 [Ulr21]。粗略地说，必须允许其状态空间是秩非常数的向量丛的场论族，其中秩的变化模仿一族 Dirac 算子的指标丛的跳跃。
\end{remark}

\subsection{来自配边框架的额外洞见}
延展 §2.3 中的想法，配边框架对带超对称的量子力学系统提供如下修饰。

\begin{example}[3.22（定向反转、KO 与 KR）]
仿照例 2.17，$\EBord_{1|1}(\mathfrak X)$ 上的定向反转通过前复合给出 $\mathfrak X$ 上（扭）场论群胚的一个 $\Z/2$-作用。取（同伦）不动点，得到无定向的 $1|1$ 维欧几里得场论（见 [HST10, Example 6.16]）。$\Z/2$-不动数据仍是状态空间上的一个实结构 (15)，并且要求时间演化算子是实的。因此，无定向理论对应于 $\KO$ 理论，见 [HST10, §6.8]。一个进一步的变体（与 Atiyah 的 $\KR$ 理论有关）以带 $\Z/2$-作用的栈 $\mathfrak X$ 为输入；这通过配边的定向反转与 $\mathfrak X$ 上的 $\Z/2$-作用两者确定 $\EBord_{1|1}(\mathfrak X)$ 上的一个 $\Z/2$-作用，进而给出场论上的 $\Z/2$-作用。
\end{example}

\begin{example}[3.23（带超联络的超向量丛）]
Dumitrescu [Dum12] 与 Stoffel [Sto19] 描述了例 2.20 与 2.21 到 $1|1$ 维欧几里得场论的推广。具体地，给定光滑流形 $X$ 上带 Quillen 超联络 $\mathbb{A}$ 的（有限秩）超向量丛 $V\to X$，Dumitrescu 构造了沿 $X$ 中 $1|1$ 维超路径的超平行移动。这给出一个函子
\begin{equation}
\mathrm{sPar}\colon\Vect^{\mathbb{A}}(X)\to\QFT^1_{1|1}(X).
\tag{38}
\end{equation}
对 $X$ 上的 gerbe $\alpha$，Stoffel 给出 (38) 的 $\alpha$-扭推广。
\end{example}

\begin{remark}[3.24]
(38) 的一个单侧逆取值于未必有限秩的向量丛上的超联络 [BE23b]。这在 [ST11] 的框架中证明了猜想 3.1 的一部分：Clifford 线性超联络的指标丛给出从场论到 K 理论所要的上闭链映射。对一族 spin$^c$ 流形，Bismut 的超联络 [Bis86] 进一步给出 (21) 中量子化映射的部分构造。要在此框架中完成猜想 3.1 的证明，只需对 Bismut 超联络构造 Dumitrescu 的超平行移动 (38) 的一个无穷秩推广，从而得到 (21) 中的“量子化”箭头。
\end{remark}
% [Tutor Note: 原文此处作 "it would suffice to to construct"，重复了 to；中文不受影响，已登记 ERRATA。]

\begin{remark}[3.25]
由于 $\EBord_{1|1}(\mathfrak X)$ 中的配边是超流形，所涉及的模空间可能很微妙，例如参见 [Wit19]。对 $\mathfrak X=\mathrm{pt}$ 上的场论，全部超几何本质上都被收纳进“超长度”参数 $(t,\theta)\in\R^{1|1}_{\ge0}$。然而，对栈 $\mathfrak X$ 上的配边，超的方面变得更加复杂。例如，对流形 $X$，$E\in\QFT^1_{1|1}(X)$ 在超圆上的取值确定一个微分形式 $Z_E\in C^\infty\bigl(\mathrm{Map}(\R^{0|1},X)\bigr)\simeq\Omega^\bullet(X)$，它与 (38) 结合起来可解释为场论族的 Chern 特征 [Han08]。更一般地，函数
% [原文 p.20]
在超栈 $\mathrm{Map}(\R^{0|1},X/\!/^{\nabla}G)$ 上给出等变 de Rham 上同调的一个模型 [Wu05, Sto13]，并引出 $1|1$-欧几里得场论的等变 Chern 特征 [BEH16]。若没有 [ST11] 的“超配边”框架，我们不知道有什么办法可以获得这些更精细的信息。
\end{remark}

\begin{example}[3.26（射影 $G$-表示）]
把超对称加到例 2.18 与 2.19 上，范畴 $\EBord_{1|1}(\mathrm{pt}/\!/^{\nabla}G)$ 把带联络的 $G$-丛加到图 4 所描绘的配边上。把超区间的半超群 $\mathrm{I}^+_{t,\theta}$ 限制到 $\theta=0$，得到一个与 (16) 所描绘的通常区间上的 $G$-丛等价的 $\EBord_{1|1}(\mathrm{pt}/\!/^{\nabla}G)$ 子范畴。因此，人们重新得到与之前相同的信息，但带有额外的增强。第一，超对称确定先前哈密顿量 $H$ 的一个平方根 $Q$。第二，超流形上 $G$-丛的联络更为复杂，涉及等变微分形式的几何。进一步讨论参见 [BEH16]。
\end{example}
